\newpage
\section{Introdução}
\label{sc:introducao}

O descasamento entre o modelo de objetos e o modelo relacional é um problema antigo da engenharia de software: linguagens orientadas a objetos organizam dados em hierarquias de classes, com identidade e associações, ao passo que bancos relacionais os representam em tabelas e chaves, sem noção de herança~\cite[p.~62]{fowler}. Os mapeadores objeto-relacionais (ORMs) surgiram para mediar essa distância, traduzindo automaticamente entre as duas representações~\cite{hibernate}. Na plataforma Java, o Hibernate~\cite{bauer-king} é um marco da categoria: lançado em 2001 como alternativa aos \emph{entity beans} da especificação EJB~2, cujo modelo de componentes impunha código verboso e descritores extensos, popularizou o mapeamento de objetos comuns a tabelas. A própria especificação seguinte, EJB~3.0 (JSR~220), nasceu com o objetivo declarado de reduzir essa complexidade do ponto de vista do desenvolvedor~\cite{jsr220}, e dela emergiu o padrão Java Persistence API (JPA), que, segundo a documentação oficial da Sun, ``incorpora as melhores ideias de tecnologias de persistência como Hibernate, TopLink e JDO''~\cite{jpa-faq}. Do grupo de especialistas da JSR~220 participou Gavin King, criador do Hibernate~\cite{jcp-king}.

Desde então, o padrão difundiu-se por praticamente todas as linguagens de uso corrente, e o ecossistema JavaScript não foi exceção. Com a adoção crescente do TypeScript, que acrescentou tipagem estática à linguagem, o mapeamento objeto-relacional ganhou um novo atrativo: declarar entidades e consultas de forma verificável pelo compilador. Formou-se daí um conjunto maduro e populoso de bibliotecas, entre elas TypeORM~\cite{typeorm}, Prisma~\cite{prisma} e Drizzle~\cite{drizzle}, cada uma com escolhas distintas de arquitetura, de modelagem e de integração com o \emph{runtime}.

Essa maturidade, porém, não dissolveu as tensões herdadas do descasamento original. Este trabalho concentra-se em três delas: a herança entre entidades, sem suporte uniforme no ecossistema, que a proposta endereça pelo mapeamento de tabela única (STI); o anti-padrão N+1, em que o acesso implícito a dados relacionados multiplica consultas e degrada o desempenho; e a manutenção do esquema na mesma linguagem da aplicação, em vez de um artefato à parte. A esses pontos soma-se a oportunidade de revisitar escolhas consolidadas da popularização da linguagem, como a padronização dos \emph{decorators} e o amadurecimento de novos \emph{runtimes}. As próximas seções detalham essa motivação e o estado da arte.
